\documentclass[11pt,a4paper]{article}
\usepackage[margin=22mm]{geometry}
\usepackage{amsmath,amssymb}
\usepackage{fontspec}
\usepackage{polyglossia}
\setdefaultlanguage{russian}
\setmainfont{DejaVu Serif}
\usepackage[unicode,hidelinks]{hyperref}

\allowdisplaybreaks[1]
\raggedright
\setlength{\parskip}{0.35em}
\newcommand{\dd}{\mathop{}\!\mathrm{d}}
\newcommand{\avg}[1]{\langle #1\rangle}
\newcommand{\ket}[1]{\lvert #1\rangle}
\newcommand{\bra}[1]{\langle #1\rvert}
\newcommand{\braket}[2]{\langle #1\mid #2\rangle}
\DeclareMathOperator{\Var}{Var}

\newcommand{\solutiontitle}[4]{%
  \hypersetup{pdftitle={Квантовая механика 1. Решения ДЗ №#1}}%
  \begin{center}
    {\small Квантовая механика 1 \enspace\textperiodcentered\enspace Семестр 1\par}
    \vspace{1ex}
    {\LARGE\bfseries Решения домашнего задания №#1\par}
    \vspace{0.5ex}
    {\large #2\par}
  \end{center}
  \noindent{\small Условия приведены в файле «#3», #4.}\par\medskip
}

\begin{document}
\solutiontitle{6}{Теорема Эренфеста и уравнения Гейзенберга}{ДЗ2026(1)\_06 (1).pdf}{с.~1}

\section{Сохранение импульса}

Для \(H=\mathbf p^2/(2m)+U(\mathbf r)\) и не зависящего явно
от времени \(p_j\) правило производной даёт
\([p_j,U]=-i\hbar\,\partial_jU\). Следовательно
\(\dot p_j=(i/\hbar)[H,p_j]=-\partial_j U\) в гейзенберговском
представлении, где правая часть берётся в тот же момент времени.

\textbf{(a)} При \(U=\alpha(x^2+y^2)\) производные
\(\partial_xU=2\alpha x,\ \partial_yU=2\alpha y,\ \partial_zU=0\).
При \(\alpha\ne0\) сохраняется \(p_z\). Проекции \(p_x,p_y\) и
импульс как вектор в общем состоянии меняются. При \(\alpha=0\)
сохраняются все три проекции.

\textbf{(b)} Пусть среднее \(p_x\) постоянно для каждого начального
состояния из общей области определения. Тогда в момент \(t=0\)
квадратичная форма \(\avg{(i/\hbar)[H,p_x]}\) равна нулю для
всех таких векторов. Вещественная и мнимая части любого
матричного элемента этой формы получаются подстановкой
\(\phi+\psi,\ \phi-\psi,\ \phi+i\psi,\ \phi-i\psi\).
Все матричные элементы обнуляются, значит \([H,p_x]=0\) на
рассматриваемой области. Постоянство в одном состоянии такого
вывода не даёт. Например, для потенциала пункта (a) чётное по \(x\)
состояние остаётся чётным, так как отражение по \(x\) коммутирует
с \(H\). Нечётный относительно отражения \(p_x\) имеет в нём
нулевое среднее при всех \(t\), хотя \([H,p_x]=2i\hbar\alpha x\).

\section{Операторы силы и скорости}

Из \([p_i,x_j]=-i\hbar\delta_{ij}\) получаем
\([H,x_j]=-i\hbar p_j/m\), так что
\(\mathbf v=\mathbf p/m\) во всех четырёх случаях.
Формула \(\mathbf F=-\nabla U\) из задачи 1 даёт
\begin{align*}
 \text{(a)}\quad &\mathbf F=(-2\alpha x,-2\alpha y,0),
 &\text{(b)}\quad &\mathbf F=(-\alpha,0,0),\\
 \text{(c)}\quad &\mathbf F=-\alpha\,\frac{\mathbf r}{r^3},
 &\text{(d)}\quad &\mathbf F=-\alpha\,\frac{\mathbf r}{r}.
\end{align*}
Для двух последних формул использовано
\(\partial_j r=x_j/r\) при \(r>0\).

\section{Эволюция двух компонент разной чётности}

Пусть \(u_0=e^{-x^2/(2\sigma^2)}/\sqrt{\sigma\sqrt\pi}\),
\(u_1=(\sqrt2 x/\sigma)u_0\), \(\sigma>0\).
Они вещественны, единичной нормы и ортогональны по чётности.
При \(\Omega=2E/\hbar\) произведение коэффициентов
\(z=c_0^*c_1=-i(\sqrt6/5)e^{-i\Omega t}\). Для эрмитова \(B\)
перекрёстный вклад в среднее есть \(2\operatorname{Re}(z B_{01})\).
С гауссовыми моментами
\(\int x^2 u_0^2\dd x=\sigma^2/2\) и
\(\int x^4 u_0^2\dd x=3\sigma^4/4\) получаем
\[
 x_{01}=\sigma/\sqrt2,\qquad
 (x^2)_{00}=\sigma^2/2,\qquad (x^2)_{11}=3\sigma^2/2 .
\]
Диагональные элементы \(x\) и внедиагональный элемент \(x^2\)
нулевые по чётности. Поэтому
\[
 \avg{x}=-\frac{2\sqrt3}{5}\sigma\sin\Omega t,\qquad
 \avg{x^2}=\frac{11}{10}\sigma^2,\qquad
 \Var(x)=\sigma^2\left(\frac{11}{10}-\frac{12}{25}\sin^2\Omega t\right).
\]
Поскольку \(Pu_0=u_0,\ Pu_1=-u_1\),
\(\avg{P}=2/5-3/5=-1/5\).
Для импульса диагональные элементы нулевые. Из
\(u_1'=(\sqrt2/\sigma)(1-x^2/\sigma^2)u_0\) следует
\((p_x)_{01}=-i\hbar/(\sqrt2\sigma)\), и потому
\[
 \avg{p_x}=-\frac{2\sqrt3\,\hbar}{5\sigma}\cos\Omega t.
\]
В формуле \(j=(\hbar/m)\operatorname{Im}(\Psi^*\partial_x\Psi)\)
диагональные слагаемые вещественны, а перекрёстная часть даёт
\((\hbar/m)\operatorname{Im}(z)(u_0 u_1'-u_1u_0')\).
Из \(u_1=(\sqrt2 x/\sigma)u_0\) разность равна
\((\sqrt2/\sigma)u_0^2\), откуда
\[
 j(x,t)=-\frac{2\sqrt3\,\hbar}{5m\sigma}\,
         \frac{e^{-x^2/\sigma^2}}{\sigma\sqrt\pi}\cos\Omega t.
\]
Интеграл по \(x\) равен \(\avg{p_x}/m\), что проверяет знак
и коэффициент потока.

\section{Гейзенберговские решения для квадратичного гамильтониана}

Уравнения \(\dot x(t)=p_x(t)/m,\quad \dot p_x(t)=-2\alpha x(t)\)
дают \(\ddot x=-(2\alpha/m)x\).
Обозначим \(x=x(0),\ p_x=p_x(0)\).
При \(\alpha>0\), \(\omega=\sqrt{2\alpha/m}\),
начальные данные \(x(0)=x,\ \dot x(0)=p_x/m\) дают
\begin{align*}
 x(t)&=x\cos\omega t+\frac{p_x}{m\omega}\sin\omega t,\\
 p_x(t)&=p_x\cos\omega t-m\omega x\sin\omega t .
\end{align*}
Импульс получен дифференцированием координаты и умножением на \(m\).
При \(\alpha=0\) имеем \(x(t)=x+p_xt/m\), \(p_x(t)=p_x\).
При \(\alpha<0\), \(\nu=\sqrt{-2\alpha/m}\),
\begin{align*}
 x(t)&=x\cosh\nu t+\frac{p_x}{m\nu}\sinh\nu t,\\
 p_x(t)&=p_x\cosh\nu t+m\nu x\sinh\nu t .
\end{align*}

\end{document}
